\section{Interaction contracts and metrics}

\paragraph{Interaction contract.}
For each benchmark item $i$, we specify a user prompt $u_i$ and a contract
$c_i$. The contract records expected response language, expected script,
required task behavior, literal spans that must be preserved, forbidden markers
or translations, and any register or locale constraints. Some contract elements
are explicit in the prompt, while others are implicit but recoverable from the
workflow. In the Spanish-English editing example above, the prompt asks for text
improvement in Spanish, but the target content is quoted in English. The
contract therefore requires English edited output, not Spanish translation.

\paragraph{First-turn global alignment.}
For model $m$ and condition $s$, we define first-turn global alignment as:
\[
\mathrm{FTGA}(i,m,s)=
\begin{cases}
1 & \text{if the first response satisfies all contract checks,}\\
0 & \text{otherwise.}
\end{cases}
\]
The checks cover language, script, preservation, task completion, and
register/locale constraints. FTGA is stricter than ordinary task success: a
response can complete the requested operation while still failing because it
uses the wrong response language, translates a quoted span, or changes a literal
file name.

\paragraph{Repair trajectory.}
If the first response fails, we append up to two standardized repair prompts
associated with the item. The repair prompts point to the violated contract
class without revealing a reference answer. The \textbf{Re-prompt Turn Tax}
(RTT) is the minimum number of repair prompts required before success:
\[
\mathrm{RTT}(i,m,s)\in\{0,1,2,3\},
\]
where $0$ denotes first-turn success and $3$ denotes unresolved after two repair
prompts. We also report Repair@1, Repair@2, and unresolved rate among initially
failed trajectories.

\paragraph{Token tax.}
Let $T_{\mathrm{total}}(i,m,s)$ be provider-reported total tokens consumed by
the trajectory through success or exhaustion of the two-repair budget, and let
$T_{0}(i,m,s)$ be the first-attempt token count. We define:
\[
\mathrm{TokenTax}(i,m,s)=\frac{T_{\mathrm{total}}(i,m,s)}{T_{0}(i,m,s)}.
\]
Token tax is a repair-burden ratio, not an API-cost estimate. A longer system
prompt can reduce repair-normalized token tax while increasing absolute tokens.
We therefore report absolute token burden separately in the analysis.
